% Companion long-record migration for Erdős #269 round 6.
% Destination label: long:actual-window-refinements
% Not frozen R. Not a blind \input. Labels may collide with the short note;
% assign destination labels as specified in R6LongRecordMigrations.md.
% Extracted from Type B edits/relocation_extracts.tex.

\section{Complements and further questions}\label{sec:open}

Theorems~\ref{res:two-prime-transcendence}
and~\ref{res:two-prime-repeated-transcendence} close both two-prime questions,
at the stronger level of transcendence.  The de-duplicated three-prime series
is not an open irrationality question: Erd\H{o}s recorded that result for any
finite set of primes in the letter of 1 January 1973
\cite[p.~335]{erdos1974letter}, without printing a proof.  The repeated
three-prime series lies outside the one-dimensional Hecke--Mahler reduction
and remains open.  At three primes the exact unresolved repeated-series
statement is best separated from the bridge, from the historical de-duplicated
assertion, and from the possible methods for proving the repeated sum.

\subsection{The actual block data and the remaining escape}

The checker already uses literal data, which we record so that the open
problems have no unspecified forcing word.  For $p\in\{2,3,5\}$ let $q,r$ be
the other two primes and put
\[
 A_p(e)=\#\{(i,j)\in\N^2:q^i r^j<p^e\},\qquad
 C_p(e)=\sum_{u=1}^{e}A_p(u).
\]
Let $I_a$ be the increasing list of internal jumps $(p,e)$ with
$p\in\{3,5\}$ and $2^a<p^e<2^{a+1}$.  For $(p,e)\in I_a$ let
\[
 \sigma_a(p,e)=\prod_{(q,f)\in I_a,\ p^e<q^f}q.
\]
Then, for $a\ge1$, the exact checker definitions are
\[
 \beta_a=2\prod_{(p,e)\in I_a}p,\qquad
 m_a^{235}=A_2(a+1)+
   \sum_{(p,e)\in I_a}(p-1)\sigma_a(p,e)
       \bigl(C_p(e)-C_2(a)\bigr).
\tag{9.1}\label{eq:actual-digit}
\]

The first formula is the four-letter radix of
Theorem~\ref{res:dyadic-alphabet}; the second is the integer implemented by the
pinned checker.  If
\[
 \nu_a=\#\{p^e:p\in\{2,3,5\},\ e\ge1,\ p^e<2^{a+1}\},
\]
its exact tested short bound is
\[
 K^{235}(B,a)=
 \left\lfloor\frac{B(\nu_a^2+10\nu_a+27)}9\right\rfloor.
\tag{9.2}\label{eq:actual-bound}
\]
The raw shell tail is the $T_a$ of Section~\ref{sec:actual-orbit}.  The
normalised state $X_a=\frac{\hgt(2^a)}2T_a$ has the mixed-radix expansion
\[
 X_a=\sum_{j\ge a}
   \frac{m_j^{235}}{\beta_a\beta_{a+1}\cdots\beta_j};
 \qquad X_{a+1}=\beta_a X_a-m_a^{235}.
\tag{9.3}\label{eq:actual-tail}
\]

The former open problem \texttt{prob:bridge269} is closed.  Its content is
Theorem~\ref{res:denominator-reduction}, with registered onset
$a_0=u+1+2v+3w$ and Lean width $90B(a+1)^2$.  The generic carry theorems of
Section~\ref{sec:escape} are now identified with the original series at that
onset.  The checker bound $K^{235}$ remains a valid finite majorant inside
the same quadratic family; it is not required for the registered equivalence.

\subsection{The intrinsic tail question}

\begin{problem}[exact nonintegrality of every reduced tail]\label{prob:tails269}
For every $B\ge1$ with $\gcd(B,30)=1$ and every $a\ge1$, prove
\[
 B X_a\notin\Z.
\tag{9.4}\label{eq:tail-nonintegrality}
\]
\end{problem}

This pointwise form is stronger-looking but cleaner than “cofinally
nonintegral”: if $B X_a$ is integral at one index, the recurrence
\[
 B X_{a+1}=\beta_a B X_a-Bm_a^{235}
\]
makes it integral at every later index.  Thus a direct solution of
Problem~\ref{prob:tails269}, joined to the proved bridge of
Theorem~\ref{res:denominator-reduction}, bypasses all residue-window
machinery.

The bounded-radix theorem gives a useful exact reduction.  Since
$2\le\beta_a\le30$, any real affine tail orbit either hits an integer or is,
cofinally often, at distance at least $1/31$ from every integer
(\lword{Erdos269/BoundedRadixTailEscape.lean}{89}{boundedRadix_zero_or_cofinal_far}{checked}).
It does not exclude the integral branch; Problem~\ref{prob:tails269} is exactly
what must do so for the actual orbit.

\subsection{A denominator-adaptive sufficient criterion}

For the literal pair $(m^{235},K^{235})$, retain the window definitions of
Section~\ref{sec:escape}:
\[
 W_{\ell,h}=\prod_{j=0}^{h-1}\beta_{\ell+j},\qquad
 F_{\ell,0}=0,\qquad
 F_{\ell,h+1}=\beta_{\ell+h}F_{\ell,h}+m_{\ell+h}^{235}.
\]

\begin{problem}[actual cofinal local-window escape]\label{prob:producer}
Prove the displayed quantifier order
\[
 \forall B\ge1\ (\gcd(B,30)=1),\ \forall a_0\ge1,\
 \exists\ell\ge a_0\ \exists h\ge1:\quad
 \operatorname{lpr}_{W_{\ell,h}}(-BF_{\ell,h})
   >K^{235}(B,\ell+h).
\tag{9.5}\label{eq:actual-escape}
\]
\end{problem}

Every $\beta_a$ is positive, so $W_{\ell,h}>0$ is automatic.  The cofinal
quantifier is present for a substantive reason: the proved bridge supplies
the reduced recurrence only after the denominator-dependent onset $a_D$,
and~\eqref{eq:actual-escape} then supplies a window beyond that onset.
Once such a window is chosen, the positive endpoint carry must equal the
canonical residue exactly
(\lword{Erdos269/RestrictedFloorSum.lean}{497}{leastPositiveResidue_windowForcing_eq_carry}{checked}),
yet it lies in the possible carry set
$\{1,\ldots,K^{235}(B,\ell+h)\}$; the strict inequality excludes that set.
This is a one-sided least-positive-residue statement, not two symmetric arcs
around zero.

Two exact countermodels rule out tempting shortcuts.  For $(W,F,B)=(6,4,1)$,
\[
 \operatorname{lpr}_{6}(-4)=2,
\]
so the canonical residue need not be coprime to $W$.  For
$(W,F,B)=(60,47,37)$,
\[
 \operatorname{lpr}_{60}(-37\cdot47)=1,
\]
so a fixed window has no denominator-independent positive residue lower
bound.  Therefore the window may genuinely depend on $B$; neither a fixed
finite computation nor a universal bounded-length guess addresses the
unbounded-denominator and cofinal-start quantifiers.  Equivalence of those
quantifiers with irrationality, in the absence of a rate, is not an
algorithmic impossibility: it names the same arithmetic in another
coordinate.  The finite checker can reject a proposed sufficient condition,
but supplies no cofinal statement.
